%% Written for PDFLaTeX
\documentclass[12pt]{article}
\usepackage[sort&compress]{natbib}
\citestyle{apalike}
\usepackage{nnonecolumn}
\usepackage{subfigure}
\usepackage{picins}
\usepackage{verbatim}
 
%% Avoid stray figures
\renewcommand{\topfraction}{0.99}
\renewcommand{\floatpagefraction}{0.99}
\renewcommand{\dbltopfraction}{0.99}
\renewcommand{\dblfloatpagefraction}{0.99}
\renewcommand{\textfraction}{.01}

%% Commands
%\newcommand{\tightenspace}[2][]{#2}      % Use if space is at a premium
\newcommand{\tightenspace}[2][]{#1}     % Use if there's plenty of space
%\renewcommand{\baselinestretch}{1.5}

\input{nnmacros.tex}

%%%%%%%%%%%%%%%%%%%%%%%%%%%%%%%%%%%%%%%%%%%%%%%%%%%%%%%%%%%%%%%%%%%%%%%%%%%%%%%
\title{{\bf Converting RGB Images\\to LMS Cone Activations}}

\author{Judah~B.~De~Paula\\\\
Department of Computer Sciences\\
The University of Texas, Austin, TX 78712-0233\\
judah@cs.utexas.edu\\\\
Report: TR-06-49 (October 2006)
}
\date{}

%%\scribbledraft



%%%%%%%%%%%%%%%%%%%%%%%%%%%%%%%%%%%%%%%%%%%%%%%%%%%%%%%%%%%%%%%%%%%%%%%%%%%%%%%
\begin{document}

\maketitle

\begin{abstract}

    This article describes a process for converting color RGB
    (Red/Green/Blue) images to LMS (Long/Medium/Short) photoreceptor
    activations.  In an RGB image each triple represents phosphor
    luminosity in a CRT (Cathode Ray Tube) screen.  In an LMS triple
    each value simulates a retina photoreceptor activation upon
    viewing the RGB triple displayed on a CRT screen.  The conversion
    process requires the phosphor emission functions for a CRT monitor
    and the sensitivity functions for the Long, Medium, and Short cone
    photoreceptors.  To calculate the final photoreceptor activity,
    transform the RGB pixels with the CRT emission function, then
    multiply the result with the photoreceptor sensitivity vector.

\end{abstract}


%%%%%%%%%%%%%%%%%%%%%%%%%%%%%%%%%%%%%%%%%%%%%%%%%%%%%%%%%%%%%%%%%%%%%%%%%%%%%%%

\section{Introduction}
\label{sec:intro}

Light enters the cornea and is focused upon the retina where the
photoreceptors activate based upon the number of photons absorbed.
There are two classes of photoreceptors called rods and cones: The
cones are used for color and daytime sight and rods are used for
low-light vision.  There are three types of cones separated into
short, medium, and long cone based upon their relative light
sensitivity functions \cite{stockman:vr99,stockman:vr00b}.  Long (L)
wavelength sensitive cones are often called red-selective cones,
medium (M) cones are called green-selective, and the short (S) are
called blue-selective.  The peak selectivity of these cells do not map
to what we would call pure red, green or blue but the names are useful
in describing their general behavior.  See the bottom plot in
Figure~\ref{fig:RGBLMS} for the normalized sensitivity curves of L, M,
and S cones as measured in humans.

Modeling a brain exposed to natural scenes requires that the
simulation must have realistic retina activation patterns.  Ideally,
the raw light from nature should be used as input to a neural network,
but natural light is often not available.  It is not possible to use a
traditional three-dimensional (RGB) digital camera to capture the
naturally occurring multi-dimensional wavelength mixtures, which would
give the most accurate information.  To reproduce a color in the
laboratory, experimenters usually use a computer monitor to present
images of natural scenes.  The RGB to LMS transformation can be used
to simulate photoreceptors viewing photographs of natural scenes.

A color image is displayed on a CRT by three types of phosphors that
are excited by electron guns.  By varying the mix of the three color
channels red, green, and blue a monitor can reproduce the majority of
colors observable by the human eye.

The goal of the CRT is to match as many colors as possible, as cheaply
as possible, while not necessarily duplicating the physical wavelength
spectrum of what the eye sees in a natural environment.  For example,
a perceived single wavelength of orange could be duplicated with a
mixture of red and yellow wavelengths.  Computer displays rely upon
this phenomenon to give the computer user an impression of natural
colors when looking at a picture.  The mixture of wavelengths from the
CRT can activate the cones in the eye in almost the same ratios as the
natural light being duplicated.  For most experiments, the
shortcomings of CRT monitors are negligible. The CRT can be calibrated
to show near lifelike images within the limits found in the camera
optics, CRT phosphors, and ambient lighting of the testing
environment.


\begin{figure}
   \centering
   \includegraphics[width=10cm,height=!]{figs/LMS/RGB_LMS.png}
   \caption{{\bf Wavelength sensitivity functions.} 
	The top plot shows RGB phosphor emission curves for a
	specific CRT monitor ($E_R^{max},E_G^{max},E_B^{max}$.)  The
	color of the curve specifies which gun was measured.  The
	bottom plot shows normalized long (red), medium (green), and
	short (blue) human cone sensitivity functions
	\cite{stockman:vr00b, stockman:vr99}.  An image designed to be
	displayed on an RGB monitor cannot be used as a photoreceptor
	activation function because it is designed for a different
	spectral curve in each channel.
	   \label{fig:RGBLMS}	}
\end{figure}


\section{Computational Modeling}

There is a problem with using RGB channels when training or testing a
neural network.  A pure wavelength stimulus, for example pure red
($R[i]=255,G[i]=0,B[i]=0$), does not stimulate only one cone type.  An
RGB picture of a natural scene has channel correlation information of
the natural environment but it is in a different form than what the
biological eye would see in nature.  It is not clear what effect the
RGB color transformation has on self-organizing neural networks, and
so it needs to be converted into LMS activation values before
training.

Figure~\ref{fig:RGBLMS} shows the difference between the light that
phosphors emit and the LMS cone sensitivity of the human eye.  Notice
how the M and L cones overlap a great deal, yet there is little
overlap in the G and R guns.

Since neuroscientists also use CRT monitors to measure brain activity,
it is reasonable to expect that the simulated and biological maps will
have similar color selectivity properties.  The next section describes
the transformation used to convert RGB images to LMS cone activations.


\section{The Transformation}

Converting from an RGB image where each triple represents phosphor
luminance, to an LMS image where each triple represents cone
stimulation, requires three data sets: 

\begin{enumerate}
\item The wavelength sensitivity function for each type of retina cone.
\item Phosphor photon emission functions for a specific computer 
      monitor.
\item RGB images of natural scenes that will be converted.
\end{enumerate}

The emission spectra of a specific CRT monitor must be measured
because each monitor has different energy emission functions.
Variations between monitors are small enough that any display with
realistic looking colors should be sufficient to use in the conversion
process.

An LMS cone activation triple can be calculated for each RGB color
pixel.  This is done by summing together the emission values for
each monitor gun at the specified RGB pixel intensity.  The
dot-product of the cone sensitivity function and the summed emission
intensity values, gives the final cone activity.

\subsubsection*{Notation}

  \begin{itemize}
    \item All uppercase letters ($L$, $M$, $S, \ldots$) represent single
	dimensional numerical row-major column vectors.  
	Subscripted vectors are	distinct: $P_1$ is different from $P_2$.
    \item Brackets denote an element within a vector.  $A[4]$ is the 
	fifth element of vector $A$.
    \item $(A)^T$ means the transpose of $A$.
    \item Lowercase letters ($\alpha$, $i$, $j, \ldots$) are scalars
	or variables.
    \item $mod$ is the remainder of integer division ($div$).
    \item All normalizations are $\infty$-norms.  The $\infty$-norm 
	of vector $A$ is the largest element of $abs(A)$.
 	
  \end{itemize}

\subsection{Step 1: Extract the R,G,B components of each image pixel}

A square RGB image named $Image$ contains normalized ($n
\times n$) triples where,

\begin{equation}
\forall i,\; i \in [0,n * n),\; 
	(R[i],G[i],B[i]) = Image[i \; div \; n][i \; mod \; n]
\end{equation}

and 

\begin{equation}
\forall i,\; i \in [0,n*n),\; 0 \leq R[i],G[i],B[i] \leq 1
\end{equation}

$R[i]$, $G[i]$, and $B[i]$ are the Red, Green, and Blue components of
the pixel $i$ within the image.

\subsection{Step 2: Calculate the monitor spectrum emission functions}

CRT monitors emit photons as a linear function relative to the input
current, however, cone luminosity sensitivity is non-linear.  Monitors
include an exponential function, usually called a gamma function, to
adjust the display.  The exponentiation of the pixel values create a
linear increase in perceptual luminosity across the bytes representing
color.

The exact gamma function may vary depending upon the monitor, but a
standard equation is: $\gamma(x) = x^{2.2}$.  Ideally, the exact
emission spectrum should be measured for each pixel value possible,
but spectral emission measurements usually only record when the
phosphor is set at its brightest value.  The top plot of
Figure~\ref{fig:RGBLMS} shows the emission vectors for when $R[i]$,
$G[i]$, and $B[i]$ are maximized.

By using the gamma function $\gamma(x) = x^{2.2}; \; x \in [0,1]$ and
the maximum spectral power distribution vectors $E_R^{max}$, $E_G^{max}$,
and $E_B^{max}$ (in Figure~\ref{fig:RGBLMS}), all possible emission
values can be calculated:

\begin{equation}
E_R(x) = \gamma(x)E_R^{max};\;
E_G(x) = \gamma(x)E_G^{max};\;
E_B(x) = \gamma(x)E_B^{max}
\end{equation}

\subsection{Step 3: Calculate total spectral power emission for each
image pixel}

To calculate the luminance power for all wavelengths coming from a CRT
pixel, the following equation is used:

\begin{equation}
\forall i,\; i \in [0,n*n),\; P_i = E_R (R[i]) + E_G (G[i]) + E_B (B[i])
\end{equation}

$E_R(\cdot)$, $E_G(\cdot)$, and $E_B(\cdot)$ each return vectors
containing the strength of the phosphor emissions for discrete
wavelengths at the specified input intensity.

\subsection{Step 4: Cone activity in the retina}

The final step is to measure the activity of a cone observing a single
CRT pixel.  This is done by taking the dot product of the spectral
sensitivity vector of the cone ($L$, $M$, or $S$, as shown in
Figure~\ref{fig:RGBLMS},) by the emission power spectrum $P_i$ of pixel
$i$ in the RGB image.

\begin{equation}
\alpha_L[i] = (P_i)^T \cdot L; \;\; \alpha_M[i] = (P_i)^T \cdot M; \;\; 
	\alpha_S[i] = (P_i)^T \cdot S 
\end{equation}

$\alpha_L[i]$, $\alpha_M[i]$, and $\alpha_S[i]$ are the scalar L, M,
and S cone activation values for image pixel $i$.


\section{Normalization}

The many stages of processing requires careful normalization of the
data.  Each subsection will discuss a different input type and what
method of normalization is required.

\subsection{CRT Spectral Power Distributions}

When the CRT is showing pure white ($R[i]=1,G[i]=1,B[i]=1$,) the peak
energy between each gun will not be the same.  Figure~\ref{fig:RGBLMS}
shows that the red gun has a higher peak than the others.  The
relationship between emission functions must be preserved since the
energy affects the photoreceptors, but it is reasonable to normalize
all three functions simultaneously to reduce rounding errors or make
the data easier to plot.  Dividing all three vectors by the wavelength
with largest energy will create a dataset that has a maximum
wavelength energy of 1.

\begin{equation}
peak = \| \;\; [\|E_R^{max}\|, \|E_G^{max}\|, \|E_B^{max}\|] \;\; \|
\label{eqn:normpeak}
\end{equation}

\begin{equation}
E_R^{norm}=\frac{E_R^{max}}{peak}, \;\;
E_G^{norm}=\frac{E_G^{max}}{peak}, \;\;
E_B^{norm}=\frac{E_B^{max}}{peak}
\label{eqn:normcrt}
\end{equation}

The area under each curve will be different, and the maximum value for
each function will be different, as they should be.

\subsection{RGB Images}

Most RGB images are already adjusted to take advantage of CRT
displays.  Some regions of the image will be dark, and some images
will contain values at or near maximum brightness.  It should not be
necessary to adjust the RGB images, but if it is required, the three
channels can be scaled using equations similar to
Equation~\ref{eqn:normpeak} and Equation~\ref{eqn:normcrt}.

Note: A corpus of training images may have biases within the separate
color channels.  For example, a picture of grass and trees may have a
Green pixel average of $0.80$ out of $1.0$, while the Blue channel
mean may be a mere $0.2$ out of $1.0$.

\subsection{Photoreceptor sensitivity functions}

Photoreceptor sensitivity functions are very difficult to measure.  It
is meaningless to have the sensitivity of one cone be greater than
another cone in the equations, unless the relationship between the
long, medium, and short cone sensitivity functions are exactly known.

This section offers three types of normalization to adjust the
photoreceptor sensitivity functions.

\subsubsection{Method 1: Independent channel normalization}
\label{sec:independ}

Assume that each cone type has an independent gain control so that
when looking at a pure white light, each cone becomes fully driven.
Each cone sensitivity function is independently normalized so that
pure white light causes an activation value of $1$.
Equation~\ref{eqn:norminde} shows the formula for this form of
normalization.

\begin{equation}
L_{norm} = \frac{L}{(L)^T \cdot P_{max}}; \;\;
M_{norm} = \frac{M}{(M)^T \cdot P_{max}}; \;\;
S_{norm} = \frac{S}{(S)^T \cdot P_{max}}
\label{eqn:norminde}
\end{equation}

This form of normalization provides the best features for training a
neural network.  Each cone type can take all possible values, and a
pure white light ($R[i]=1,G[i]=1,B[i]=1$) generates maximum cone
activation ($L[i]=1,M[i]=1,S[i]=1$.)


\subsubsection{Method 2: Preserve channel correlation}  

This method preserves the relationship in sensitivity between the
three channels in the original data.  The cone sensitivity functions
are normalized similar to Equation~\ref{eqn:normcrt}, but the
normalization term $peak$ is calculated using
Equation~\ref{eqn:normlmspeak}.  Looking at a pure white display
(Equation~\ref{eqn:purewhite},) one of the three cones will have the
highest activity, possibly beyond the required limit of $1$.  The
highest cone activity is then used as the normalization term $peak$.

\begin{equation}
P_{max} = E_R^{max} + E_G^{max} + E_B^{max}
\label{eqn:purewhite}
\end{equation}

\begin{equation}
peak = \| \;\; [ (L)^T \cdot P_{max}, 
	(M)^T \cdot P_{max}, 
	(S)^T \cdot P_{max}] \;\; \|
\label{eqn:normlmspeak}
\end{equation}

After normalization, a pure white display will cause the most
sensitive cone to fully activate to a value of $1$, while the other
two channels will peak out at some value less than $1$.  This assumes
that there is no channel independent gain control, unlike the first
method.


\subsubsection{Method 3: Find the relationship}

Find a dataset that has an exact peak-sensitivity correlation between
the three cone types.  This also requires an understanding of the
behavior of the photoreceptors under different lighting conditions not
included in this model.


\section{Conclusion}

Upon completion of the RGB to LMS conversion, each triple of
information contains simulated retinal cone activations as if the
retina contained three overlapping regions of L,M, and S cones.



\section{Appendix: Python Conversion Code}

Python code is available from the author that performs the above
transformation.

% \begin{small}
% \verbatiminput{../rgbtolms}
% \end{small}


%%%%%%%%%%%%%%%%%%%%%%%%%%%%%%%%%%%%%%%%%%%%%%%%%%%%%%%%%%%%%%%%%%%%%%%%%%%%%%%
\bibliographystyle{nnapalike}
%%\bibliographystyle{achicago}
\bibliography{nnstrings,nn,judah}
\end{document}
%%%%%%%%%%%%%%%%%%%%%%%%%%%%%%%%%%%%%%%%%%%%%%%%%%%%%%%%%%%%%%%%%%%%%%%%%%%%%%%
